\section{Objetivos}

\subsection{Objetivo General}
Disponer de un modelo predictivo de alertas tempranas, entrenado con datos de SECOP II para contratos de infraestructura financiados con recursos del Sistema General de Regalías (SGR), que clasifique el riesgo de incumplimiento o irregularidad desde la etapa de adjudicación y alcance un desempeño mínimo de AUC-ROC $\geq 0.75$ y \textit{recall} $\geq 0.70$ en validación cruzada estratificada de 5 pliegues.

\subsection{Objetivos Específicos}
\begin{enumerate}
    \item Definir un conjunto validado de variables contractuales y financieras (por ejemplo: anticipo, modalidad de contratación, número de oferentes, valor y ubicación geográfica), acompañado de un ranking de importancia para la predicción del riesgo.

    \item Consolidar una base de datos analítica integrada y depurada de SECOP II para contratos SGR, con diccionario de datos, trazabilidad del preprocesamiento y un nivel de completitud de al menos 85\% en las variables críticas definidas por el proyecto.

    \item Obtener tres modelos supervisados entrenados sobre el mismo conjunto de datos (Regresión Logística, Random Forest y XGBoost), listos para su comparación bajo un protocolo único de evaluación.

    \item Generar una evaluación comparativa reproducible de los modelos, con validación cruzada estratificada ($k=5$) y reporte de AUC-ROC, precisión, \textit{recall} y F1-score.

    \item Seleccionar un modelo final mediante criterios explícitos de desempeño (máximo AUC-ROC y \textit{recall} mínimo de 0.70), junto con su configuración de hiperparámetros documentada.

    %\item Entregar una salida operativa de priorización de riesgo (índice de riesgo y listado de contratos priorizados) y una guía de uso para entidades de control y veeduría ciudadana.
\end{enumerate}